%%% FIELD prop-finite-insulin-demonstration-proof
Put \(\mathcal G=\{50,100,150,200,250\}\), \(\mathcal D=\{0,78\}\), and \(\mathcal E=\{0,31,850\}\).  By \(\mathrm{def{:}insulin\text{-}grid}\), a joint state is an element of
\[
 \mathcal G\times\mathcal D^2\times\mathcal E^2\times\{0,1\}.
\]
Consequently its cardinality is
\[
 5\cdot 2^2\cdot 3^2\cdot 2=360.
\]

We first write the transition matrix explicitly.  For
\(s=(g,d_1,d_2,x_1,x_2,a_1)\) and \(a\in\{0,1\}\), define
\[
 m(s,a)=10+0.9g+0.1d_1+0.1d_2-0.01x_1-0.01x_2-2a-4a_1.
\]
Let \(\Phi\) be the standard normal distribution function, let
\(c_0=-\infty\), \(c_1=75\), \(c_2=125\), \(c_3=175\),
\(c_4=225\), and \(c_5=+\infty\), and enumerate the elements of
\(\mathcal G\) increasingly as \(g_0,\ldots,g_4\).  The probability
that the rounded next glucose is \(g_i\), conditional on the
nonrefresh branch, \(s\), and \(a\), is
\[
 q_i(s,a)=\Phi\!\left(\frac{c_{i+1}-m(s,a)}5\right)
             -\Phi\!\left(\frac{c_i-m(s,a)}5\right),
 \qquad 0\le i\le4,
\]
where \(\Phi(-\infty)=0\) and \(\Phi(+\infty)=1\).  This is exactly
the clipping-and-nearest-grid rule in \(\mathrm{def{:}insulin\text{-}grid}\),
since its four rounding cut points are \(75,125,175,225\); ties have
normal probability zero.  Write
\[
 \pi_D(0)=\frac45,\quad \pi_D(78)=\frac15,
 \qquad
 \pi_E(0)=\pi_E(31)=\frac25,\quad \pi_E(850)=\frac15.
\]
Thus the nonrefresh action-specific kernel \(B_a\) is specified by
\[
 B_a\bigl(s,(g_i,d_0,d_1,x_0,x_1,a)\bigr)
   =q_i(s,a)\pi_D(d_0)\pi_E(x_0),
\]
and is zero at every other next state.  The two nonrefresh
policy-induced kernels are therefore
\[
 Q_b=\frac7{10}B_0+\frac3{10}B_1,
 \qquad
 Q_e(s,\cdot)=B_{\mathbf 1\{g\ge125\}}(s,\cdot).
\]
If \(u\) denotes the uniform distribution on the \(360\) states, the
full policy-induced kernels are
\[
 P_p(s,\cdot)=\frac12u+\frac12Q_p(s,\cdot),
 \qquad p\in\{b,e\}.
\]

For any two state distributions \(\nu,\nu'\), the signed measure
\(\eta=\nu-\nu'\) has total mass zero, whence the common refresh term
cancels and
\[
 \eta P_p=\frac12\eta Q_p.
\]
A Markov kernel is a contraction in total variation: indeed, for
every function \(f\) with \(0\le f\le1\), also \(0\le Q_pf\le1\), so
\(\sup_{0\le f\le1}|\eta Q_pf|\le\sup_{0\le h\le1}|\eta h|\).
It follows that
\[
 \|\nu P_p-\nu'P_p\|_{\mathrm{TV}}
 \le\frac12\|\nu-\nu'\|_{\mathrm{TV}},
 \qquad p\in\{b,e\}.
\]
Hence both Dobrushin coefficients are at most \(1/2\).

The behavior probabilities of actions \(0\) and \(1\) are \(7/10\)
and \(3/10\).  If \(g\ge125\), the only action having positive target
probability is action \(1\), and its target-to-behavior ratio is
\(1/(3/10)=10/3\).  If \(g<125\), the only such action is action
\(0\), and its ratio is \(1/(7/10)=10/7\).  Thus the maximum one-step
ratio is
\[
 L=\max\left\{\frac{10}{3},\frac{10}{7}\right\}=\frac{10}{3}.
\]

The refresh representation also gives, for every pair of states
\(s,s'\),
\[
 P_b(s,s')\ge \frac12u(s')=\frac1{720}.
\]
Stationarity and positivity of \(d_b\) then imply
\[
 d_b(s')=\sum_s d_b(s)P_b(s,s')\ge\frac1{720}\sum_s d_b(s)
 =\frac1{720}.
\]
Since \(d_e(s')\le1\), division by the strictly positive quantity
\(d_b(s')\) yields
\[
 d_e(s')\le720d_b(s')
\]
for every joint state \(s'\).  This is the claimed analytic latent
stationary-overlap certificate.

It remains to certify that immediate weighting is biased.  The reward
is a deterministic function \(r\) of the current glucose, with
\[
 r(50)=-1,\quad r(100)=0,\quad r(150)=-\frac13,\quad
 r(200)=r(250)=-\frac23.
\]
Both behavior action probabilities are positive.  Conditional
behavior randomization and the fact that \(r\) does not depend on the
current action give
\[
 \begin{aligned}
 E_b\!\left[Y_t\frac{e(A_t\mid X_t)}{b(A_t\mid X_t)}\right]
 &=\sum_s d_b(s)r(s)\sum_{a=0}^1
       b(a\mid x)\frac{e(a\mid x)}{b(a\mid x)}\\
 &=\sum_s d_b(s)r(s).
 \end{aligned}
\]
On the other hand, \(\mathrm{def{:}target\text{-}value}\) and the same
action-independence of the current reward give
\[
 \theta(K,e)=\sum_s d_e(s)r(s).
\]
Therefore
\[
 \Delta_{\mathrm{imm}}=d_br-d_er.
\]

Here is a reproducible rational interval certificate for this last
difference.  From \(P_p=\tfrac12u+\tfrac12Q_p\), define
\[
 \bar d_p:=\sum_{j=0}^{\infty}2^{-(j+1)}uQ_p^j.
\]
The series is an absolutely convergent probability vector, and direct
multiplication gives
\[
 \bar d_pP_p=\frac12u+\frac12\bar d_pQ_p
 =\frac12u+\sum_{j=1}^{\infty}2^{-(j+1)}uQ_p^j=\bar d_p.
\]
The total-variation contraction just proved makes the stationary law
unique, so \(\bar d_p=d_p\).

All normal probabilities used below were enclosed using rational
arithmetic.  For completeness, the enclosure is now described.  For
\(0\le x\le8\), termwise integration of the exponential power series
gives
\[
 \Phi(x)=\frac12+\frac1{\sqrt{2\pi}}
 \sum_{k=0}^{\infty}\frac{(-1)^kx^{2k+1}}{(2k+1)2^kk!},
 \qquad \Phi(-x)=1-\Phi(x).
\]
The constant was enclosed without floating-point evaluation.  The
tangent addition formula gives
\[
 \tan\!\left(4\arctan\frac15-\arctan\frac1{239}\right)=1,
\]
because \(\tan(2\arctan(1/5))=5/12\),
\(\tan(4\arctan(1/5))=120/119\), and substituting
\(1/239\) in the subtraction formula gives exactly one.  The angle
lies in \((0,\pi/2)\), and hence
\[
 \pi=16\arctan\frac15-4\arctan\frac1{239}.
\]
For \(0<x<1\), integrating the finite geometric identity for
\((1+x^2)^{-1}\) shows that the truncation
\[
 A_N(x)=\sum_{k=0}^N\frac{(-1)^kx^{2k+1}}{2k+1}
\]
has an error of the sign of its first omitted term and magnitude at
most \(x^{2N+3}/(2N+3)\).  Applying this with \(N=14\) at \(1/5\)
and \(N=4\) at \(1/239\), followed only by integer multiplication and
comparison, gives
\[
 \frac{3989422804014326}{10^{16}}
 <\frac1{\sqrt{2\pi}}<
 \frac{3989422804014328}{10^{16}}.
\]
For the integrated exponential series, the ratio of consecutive
absolute terms is at most \(x^2/(2(k+1))\); hence its tail is
alternating and decreasing after \(k=32\).  Truncation after \(k=118\)
therefore brackets the integral between that partial sum and the same
sum minus its next term.  Uniformly for \(x\le8\), that next term is at
most
\[
 \frac{8^{239}}{239\,2^{119}119!}<8\times10^{-20}.
\]
For \(x>8\), integration of
\(e^{-t^2/2}\le(t/x)e^{-t^2/2}\) for \(t\ge x\) gives the Mills bound
\[
 1-\Phi(x)\le\frac{e^{-x^2/2}}{x\sqrt{2\pi}}.
\]
The exact rational comparison
\[
 \sum_{k=0}^{40}\frac{32^k}{k!}>
 \frac{500000000000000}{7}
\]
and the preceding upper bound \(1/\sqrt{2\pi}<2/5\) imply
\[
 1-\Phi(x)<7\times10^{-16},\qquad x>8.
\]

There are \(1422\) distinct rational arguments
\[
 \mathcal Z=\left\{\frac{c-m(s,a)}5:
 s\in\mathcal G\times\mathcal D^2\times\mathcal E^2\times\{0,1\},
 \ a\in\{0,1\},\ c\in\{75,125,175,225\}\right\},
\]
and their range is \([-878/25,193/5]\).  Applying the preceding
enclosures to this finite set, exact rational comparison places every
interval strictly inside one rounding cell for the grid
\(10^{-12}\mathbb Z\); the smallest clearance from a cell boundary is
greater than \(7\times10^{-4}\) after multiplication by \(10^{12}\).
Let \(\widetilde\Phi(z)\) be the resulting multiple of \(10^{-12}\).
Then, simultaneously for all \(z\in\mathcal Z\),
\[
 |\Phi(z)-\widetilde\Phi(z)|\le\varepsilon,
 \qquad \varepsilon:=5\times10^{-13}.
\]

Replace \(\Phi\) by \(\widetilde\Phi\) in the displayed construction
of \(B_a,Q_b,Q_e\), obtaining stochastic matrices
\(\widetilde Q_b,\widetilde Q_e\).  Monotonicity of rounding and
telescoping of the five glucose probabilities make their entries
nonnegative and their row sums one.  Each action-specific row has
\(\ell^1\) error at most
\[
 \varepsilon+3(2\varepsilon)+\varepsilon=8\varepsilon;
\]
averaging over an action cannot increase this bound.  Consequently
\[
 \sup_s\|Q_p(s,\cdot)-\widetilde Q_p(s,\cdot)\|_{\mathrm{TV}}
 \le4\varepsilon,\qquad p\in\{b,e\}.
\]

The following exact-integer recurrence records the remaining finite
calculation.  Order the states lexicographically, put
\(R=10^{12}\), and put \(D_Q=250R\).  The matrices
\[
 M_p:=D_Q\widetilde Q_p
\]
have nonnegative integer entries and every row sum is \(D_Q\): the
nonrefresh glucose probabilities have denominator \(R\), the diet and
activity probabilities together have denominator \(25\), and the
behavior mixture has denominator \(10\); the target rows are put over
the same denominator by multiplication by \(10\).  Starting with
\(n_{p,0,i}=1\), iterate
\[
 n_{p,j+1,k}=\sum_{i=0}^{359}n_{p,j,i}M_p(i,k).
\]
Thus \(u\widetilde Q_p^j\) has coordinate
\(n_{p,j,i}/(360D_Q^j)\).  With \(\bar r_i=3r(s_i)\in\{-3,-2,-1,0\}\),
define the entirely rational quantity
\[
 \widetilde V_{p,30}
 =\sum_{j=0}^{29}\frac1{2^{j+1}}
   \frac{\sum_{i=0}^{359}n_{p,j,i}\bar r_i}{3\cdot360D_Q^j}.
\]
Exact integer evaluation of this displayed recurrence gives
\[
 -\frac{5906084184470}{10^{15}}
 <\widetilde V_{b,30}-\widetilde V_{e,30}
 <-\frac{5906084184469}{10^{15}}.
\]
This calculation involves no floating-point operation: the normal
enclosures select integers, and the recurrence thereafter consists
only of integer additions and multiplications.

Finally, since \(r\in[-1,0]\), its oscillation is one.  A telescoping
comparison of \(Q_p^j\) and \(\widetilde Q_p^j\), using the last
total-variation row bound at each of the \(j\) positions, yields
\[
 |uQ_p^jr-u\widetilde Q_p^jr|\le4j\varepsilon.
\]
Also, the omitted geometric tail has absolute value at most
\(2^{-30}\).  Since \(\sum_{j\ge0}j2^{-(j+1)}=1\), the two policies
together satisfy
\[
 \left|\Delta_{\mathrm{imm}}-
   (\widetilde V_{b,30}-\widetilde V_{e,30})\right|
 \le8\varepsilon+2^{1-30}<1.867\times10^{-9}.
\]
Combining this rigorous error bound with the exact rational interval
above gives the sharper certificate
\[
 -0.005906087<\Delta_{\mathrm{imm}}<-0.005906082.
\]
In particular,
\[
 -0.0060<\Delta_{\mathrm{imm}}<-0.0058,
\]
which proves the asserted nonzero immediate-weighting bias together
with all the structural claims.
